\subsection{Scaling Graph Neural Networks to Large Graphs}
\label{sec:relwork:scaling}

\subsubsection{Sampling-Based Approaches}
\label{sec:relwork:scaling:sampling}
The oldest answer to a graph that does not fit is to train on samples of it: neighbour
sampling caps the fanout of each aggregation \citep{hamilton2017graphsage}, and subgraph
sampling trains on induced pieces. Sampling is a per-epoch stochastic operation built for
node-level targets, and both properties disqualify it here. A graph-level regression
needs every node of a graph in one forward pass (\ref{sec:prelim:gnn:readout}), and the
reduction under study must be a deterministic, cacheable artifact computed once offline,
so that its cost amortises across runs and its output can be inspected as a graph in its
own right.

\subsubsection{Partitioning and Distributed Training}
\label{sec:relwork:scaling:distributed}
The standard industrial answer is to partition the graph across devices and exchange
messages at the boundaries, or forgo them there. This is the setting in which cutting
edges is forced rather than chosen, and it is why partitioning is carried as a reduction
family here (\ref{sec:prelim:reduction:partitioning}) despite removing no nodes: the
accuracy cost of severed spanning edges is exactly what a distributed practitioner pays,
measured on one device.
% TODO(citation): add a distributed GNN training reference.

\subsubsection{Memory-Efficient Training Techniques}
\label{sec:relwork:scaling:memory}
Gradient checkpointing, mixed precision and gradient accumulation, defined
in~\ref{sec:prelim:gnn:cost}, are used in this thesis rather than studied by it. They
must still be named, because a reader will ask why input reduction is needed when
checkpointing exists. The answer is composition: each of these techniques trades
compute, precision or latency for memory at a fixed input size, reduction changes the
input size itself, and the two stack. The baseline configurations depend on the first
fact, DeepGate4 being runnable only under checkpointing
(\ref{sec:method:architecture:baselines}), and the contribution depends on the second.
